\documentclass[a4paper,11pt]{article}
\usepackage[utf8]{inputenc}
\usepackage[T1]{fontenc}
\usepackage[french]{babel}
\usepackage{lmodern}
\usepackage{amsmath,amssymb,amsthm}
\usepackage{mathrsfs}
\usepackage{enumitem}
\usepackage{multicol}

\setlist[enumerate]{label=\textbf{\textcolor{Couleur8}{\arabic*.}}, left=1em}

% Si vous voulez aussi modifier les listes enumerate imbriquées :
\setlist[enumerate,2]{label=\textcolor{Couleur8}{\alph*)}}
\setlist[enumerate,3]{label=\textcolor{Couleur8}{\roman*)}}
\setlist[enumerate,4]{label=\textcolor{Couleur8}{\Alph*)}}
\usepackage{graphicx}
\usepackage{xcolor}
\usepackage{tcolorbox}
\usepackage{geometry}
\usepackage{fancyhdr}
\usepackage{titlesec}
\usepackage{cancel}
\usepackage{ifthen}
\usepackage{tikz}
\usetikzlibrary{arrows.meta}
\usepackage{tkz-tab}
\usepackage{eurosym}

% OPTION PROF/ÉLÈVE
% Décommentez UNE des deux lignes suivantes :
\newboolean{prof}\setboolean{prof}{true}   % Version professeur (avec corrections des devoirs)
\newboolean{prof}\setboolean{prof}{true}  % Version élève (sans corrections des devoirs)

\definecolor{Couleur1}{HTML}{4A5D7A} %Th
\definecolor{Couleur2}{HTML}{A5B5D6} %Prop
\definecolor{Couleur3}{HTML}{A8C68A} %Méthode
\definecolor{Couleur6}{HTML}{8A9CC1} %Def
\definecolor{Couleur7}{HTML}{E6B459} %Rq Voc
\definecolor{Couleur8}{HTML}{D36740} %Titres
\definecolor{correction}{HTML}{D36740} % Couleur pour les corrections
\definecolor{coopmathscorr}{HTML}{F15929} % Couleur corrections coopmaths

% Configuration des marges
\geometry{
	top=2cm,
	bottom=2cm,
	inner=1.5cm,
	outer=1.5cm,
}

% Police
\renewcommand{\familydefault}{\sfdefault}

% Compteurs
\newcounter{seancenum}
\newcounter{seancenumD}
\setcounter{seancenum}{1}
\setcounter{seancenumD}{1}

% Configuration des en-têtes et pieds de page
\pagestyle{fancy}
\fancyhf{}
\ifthenelse{\boolean{prof}}{
    \fancyhead[L]{\textcolor{Couleur8}{\textbf{Corrections Automatismes - Classe de seconde - Version Professeur}}}
}{
    \fancyhead[L]{\textcolor{Couleur8}{\textbf{Corrections Devoirs - Séance 27 - Classe de seconde}}}
}
\fancyhead[R]{\textcolor{Couleur8}{\textbf{2025-2026}}}
\fancyfoot[L]{\textcolor{Couleur8}{Mme Bertail et Mme Pinto}}
\fancyfoot[R]{\textcolor{Couleur8}{\thepage}}
\renewcommand{\headrulewidth}{1pt}
\renewcommand{\headrule}{\hbox to\headwidth{\color{Couleur8}\leaders\hrule height \headrulewidth\hfill}}

% Environnements personnalisés
\newtcolorbox{seance}[1]{
    colback=Couleur6!10,
    colframe=Couleur1!65,
    fonttitle=\bfseries\large,
    title=Correction Séance \theseancenum #1,
    left=5mm,
    right=5mm,
    top=3mm,
    bottom=3mm,
    arc=0mm,
    after=\stepcounter{seancenum}
}

\newtcolorbox{seanceD}[1]{
    colback=Couleur6!5,
    colframe=Couleur6!45,
    fonttitle=\bfseries\large,
    title=Correction Devoirs - Séance \theseancenumD #1,
    left=5mm,
    right=5mm,
    top=3mm,
    bottom=3mm,
    arc=0mm,
    after=\stepcounter{seancenumD}
}

\begin{document}

\setcounter{seancenumD}{27}
\newpage
\begin{seanceD}{} %27

\begin{enumerate}
\item On met l'expression au même dénominateur : \\$\begin{aligned}
        \dfrac{1}{2}-\dfrac{2x+5}{x}&=\dfrac{x-2\times \left(2x+5\right)}{2x}\\
        &=\dfrac{x -4x -10}{2x}\\
        &=\dfrac{-3x -10}{2x}\\
     \end{aligned}$\\$\phantom{\dfrac{1}{2}-\dfrac{2x+5}{x}}=-\dfrac{3x +10}{2x}$\\La bonne réponse est la réponse {\bfseries \color[HTML]{F15929}B}.

\item On applique la propriété du quotient des puissances d'un réel : \\
      Soit $n$ et $p$ deux entiers et $a$ un réel :  $\dfrac{a^n}{a^p}=a^{n-p}$\\
      $\begin{aligned} \dfrac{a^{n}}{a^{n^{5}}}&=a^{n-n^{5}}\\
      &=a^{-n(-1+n^{4})}\\
      &={\color[HTML]{F15929}\boldsymbol{a^{-n(n^{4}-1)}}}
      \end{aligned}$

\medskip
La bonne réponse est la réponse {\bfseries \color[HTML]{F15929}D}.

\item On reconnaît une différence de deux carrés $A^2-B^2$ avec $A=7x+4$ et $B=5x-1$, qui se factorise en $(A-B)(A+B)$ .\\
    $\begin{aligned}
    \left(7x+4\right)^2-\left(5x-1\right)^2&=\left[\left(7x+4\right)-\left(5x-1\right)\right]\left[\left(7x+4\right)+\left(5x-1\right)\right]\\
    &=\left(2x+5\right)\left(12x+3\right)\\
    &={\color[HTML]{F15929}\boldsymbol{24x^{2}+66x+15}}.
    \end{aligned}$\\La bonne réponse est la réponse {\bfseries \color[HTML]{F15929}B}.

\item On a :\\
          $\begin{aligned}
          f(-5)&=3- 2\times (-5)^2\\
          &={\color[HTML]{F15929}\boldsymbol{-47}}
          \end{aligned}$\\{\color[HTML]{216D9A} Mentalement : \\
    On commence par "calculer" le carré de $-5$ :  $(-5)^2=25$.\\
    Puis on multiplie le résultat par $2$ : $2\times 25=50$.\\
    On calcule alors : $3-50=-47$.}\\La bonne réponse est la réponse {\bfseries \color[HTML]{F15929}C}.

\item La droite $(D)$ est strictement au-dessus de la droite $(D')$ lorsque $5x+5>8x+1$.\\
     $\begin{aligned}
     5x+5&>8x+1\\
     -3x&>-4\\
     x&<\dfrac{4}{3}
     \end{aligned}$\\
     La droite $(D)$ est donc strictement au-dessus de la droite $(D')$ sur : 
     ${\color[HTML]{F15929}\boldsymbol{\left]-\infty\,;\,\dfrac{4}{3}\right[}}$.\\La bonne réponse est la réponse {\bfseries \color[HTML]{F15929}B}.

\item On note $A$ l'événement \og{} réussir au moins une fois un tir \fg{} et $B$ l'événement \og{} rater tous les tirs \fg{}.\\
Les événements $A$ et $B$ sont contraires.\\
Donc : $P(B) = 1 - P(A)$.\\
$P(B) = 1 - 0{,}992 = 0{,}008$\\La bonne réponse est la réponse {\bfseries \color[HTML]{F15929}D}.

\end{enumerate}

\end{seanceD}

\end{document}